% ============================================================
%  Observation equations
% ============================================================

This section gives the \emph{undifferenced} observation model about
which the filter linearises at each epoch.  The PPP-RTK
parameterisation is uncombined: the filter consumes per-signal
pseudorange and carrier phase directly and forms no ionosphere-free or
wide-lane combination as a preprocessing step.  The combinations
(\textsc{IF}, \textsc{MW}, \textsc{GF}) are computed for diagnostics and
slip detection only (\dcref{sec:slip}{the cycle-slip document};
\texttt{slip/MwSlipDetector}, \texttt{slip/GraphicSlipDetector}) and do
not enter the measurement update.  The loader
(\texttt{GnssConfig::loadProcessingSection\_}) forces
\texttt{uncombined\_enabled=true} for the whole PPP-RTK family (user,
network producer, network consumer) and for the baseline
double-difference modes, which process each signal independently
(\Cref{sec:dd}).  Plain PPP mode keeps the flag at its default
\texttt{false} and builds ionosphere-free combined observables (keys
such as \texttt{C1W\_C5X}); these feed
\texttt{MeasurementModel::buildJacobianRow} and the update.  The rest of
this section describes the uncombined case.

\subsection{Code (pseudorange) observation}

For receiver $\rx$ tracking satellite $\sat$ on signal $\sigl$, the
pseudorange is
\begin{empheq}[box=\fbox]{equation}
\code^\sat_{\rx,\sigl}
\;=\;
\geomrange^\sat_\rx
\;+\; \sLight(\mathrm{d}t_\rx - \mathrm{d}t^\sat)
\;+\; \sLight\,\delta t^{(g(\sat))}_\rx
\;+\; T^\sat_\rx
\;+\; \mu_\sigl\, I^\sat_{\rx,1}
\;+\; B_{\rx,\sigl} - B^\sat_\sigl
\;+\; \varepsilon^\code_{\rx,\sat,\sigl}.
\label{eq:obs-code}
\end{empheq}

\subsection{Carrier-phase observation}

The carrier-phase observation, scaled to metres by $\lambda_\sigl$, is
\begin{empheq}[box=\fbox]{equation}
\phaseo^\sat_{\rx,\sigl}
\;=\;
\geomrange^\sat_\rx
\;+\; \sLight(\mathrm{d}t_\rx - \mathrm{d}t^\sat)
\;+\; \sLight\,\delta t^{(g(\sat))}_\rx
\;+\; T^\sat_\rx
\;-\; \mu_\sigl\, I^\sat_{\rx,1}
\;+\; \lambda_\sigl\, N^\sat_{\rx,\sigl}
\;+\; b_{\rx,\sigl} - b^\sat_\sigl
\;+\; \varepsilon^\phaseo_{\rx,\sat,\sigl}.
\label{eq:obs-phase}
\end{empheq}

The two equations differ structurally in the sign of the ionospheric
term: $+\mu_\sigl I$ on code, $-\mu_\sigl I$ on phase.  Further
remarks:

\begin{enumerate}[leftmargin=*, itemsep=0pt]
\item The phase ambiguity $N^\sat_{\rx,\sigl}$ is an integer in
      principle, but the receiver and satellite phase biases
      $b_{\rx,\sigl} - b^\sat_\sigl$ contaminate it.  A \emph{float}
      solution absorbs these biases into $N^\sat_{\rx,\sigl}$, which
      leaves a non-integer state with no structural advantage over the
      code observation.  Integer-valued ambiguities therefore require
      either external Fractional Cycle Bias (FCB) products or
      differencing (\Cref{sec:diff}), which cancels the common-mode
      bias terms.
\item The clock product $\sLight\,\mathrm{d}t^\sat$ comes from the
      analysis-centre CLK file (post-processing) or from the broadcast
      clock refined by a streamed SSR correction (real time,
      \dcref{sec:ssr}{the SSR-corrections document}).  The filter treats
      it as deterministic.  Its error enters the observation: the part
      common to all satellites is absorbed by the receiver clock state,
      and the satellite-specific part remains in the residuals.
\item The system time offset $\sLight\,\delta t^{(g(\sat))}_\rx$ is
      zero by construction when $g(\sat)$ is the reference
      constellation, which is GPS (the state manager keys no offset for
      GPS).  With the receiver clock $\sLight\,\mathrm{d}t_\rx$, a
      multi-constellation run therefore carries $|\sysset|$ clock-like
      states in total.
\item The slant tropospheric delay is the hydrostatic
      zenith delay mapped by the dry mapping function plus the wet
      zenith delay mapped by the wet mapping function
      (\texttt{PppObservation} fields \texttt{mf\_d} and \texttt{mf\_w}).
\item The equations omit the deterministic model corrections applied to
      the computed range: satellite and receiver antenna phase-centre
      offsets and variations, the relativistic terms, site displacement,
      and, on phase rows only, the carrier-phase wind-up.  The wind-up
      follows the \texttt{nominal-wu93} sign convention.
\end{enumerate}

\subsection{Geometric range and linearisation}\label{sec:obs-geometry}

The geometric range $\geomrange^\sat_\rx = \|\satpos - \rxpos_\rx\|_2$
enters \Cref{eq:obs-code,eq:obs-phase} non-linearly through the
receiver position $\rxpos_\rx$.  The filter linearises about the
current estimate $\hat{\rxpos}_\rx$:
\begin{equation}
\geomrange^\sat_\rx
\;=\;
\hat{\geomrange}^\sat_\rx
\;-\; (\los^\sat_{\hat{\rx}})^{\Tr} (\rxpos_\rx - \hat{\rxpos}_\rx)
\;+\; \mathcal{O}\!\left(\|\rxpos_\rx - \hat{\rxpos}_\rx\|^2\right),
\label{eq:geom-linear}
\end{equation}
with the unit line-of-sight
\begin{equation}
\los^\sat_{\hat{\rx}}
\;:=\;
\frac{\satpos - \hat{\rxpos}_\rx}{\|\satpos - \hat{\rxpos}_\rx\|_2}.
\label{eq:los}
\end{equation}
The position columns of the Jacobian row $\jacH_{(\sat,\sigl)}$ are
therefore $-(\los^\sat_{\hat{\rx}})^{\Tr}$.  One linearisation per epoch
suffices for sub-metre residual targets.  For large initial position
offsets the engine can relinearise within the epoch
(\texttt{processing.relinearize\_geometry}).  The flag defaults to off in
code; the PPP, PPP-RTK, baseline and real-time configs in
\texttt{config/} enable it, and the six network-producer configs leave
it off.

The satellite position $\satpos$ is a function of the transmission
time.  \texttt{Geometry::computeEmissionState} obtains the emission time
in one step from the \emph{pseudorange} rather than by iterating against
the geometric range, which depends on $\rxpos_\rx$.  The time of flight
$\tau = \code / \sLight$ and the satellite clock are subtracted from the
receiver epoch.  The receiver clock bias $\mathrm{d}t_\rx$ enters the receiver epoch as
$+\mathrm{d}t_\rx$ and the pseudorange as $+\sLight\,\mathrm{d}t_\rx$,
so it cancels exactly.  The remaining error $(I+T+\varepsilon)/\sLight
\approx 30$--$60$~ns displaces the satellite by less than a millimetre,
below any PPP threshold.  The seed pseudorange is the ionosphere-free
combination of the clock-reference signals when available
(\texttt{Geometry::extractPseudorange}), so the ionospheric part of that
error applies only to the single-frequency fallback.  A range-based fixed
point is avoided on purpose: iterating against $\geomrange^\sat_\rx$
would reintroduce an uncancelled $\mathrm{d}t_\rx$ and drive a
positive-feedback loop through the receiver-clock state.

The one second pass is the \texttt{model\_corrections.refine\_emission\_time}
block in \texttt{Geometry::computeSatelliteGeometry} (default on).  It
re-runs \texttt{computeEmissionState} on a cleaned pseudorange: the seed
minus the slant hydrostatic delay and the IF-combined ATX projection,
plus the periodic relativistic clock term.  This refines the
pseudorange; it is not a range-versus-position iteration.

\subsection{Stochastic model}

The observation noise covariance of an epoch is diagonal, with entries
\begin{equation}
\sigma^2_{\rx,\sat,\sigl}
\;=\;
\sigma^2_{0,\text{type}(\sigl)}\;\cdot\; w(\theta^\sat_\rx)^{-1}\;\cdot\;
\xi(\text{C/N}_0\text{ at } (\rx,\sat,\sigl)),
\label{eq:noise-model}
\end{equation}
The three factors are:
\begin{itemize}[leftmargin=*, itemsep=0pt]
\item a per-type baseline ($\sigma_{0,\text{code}}$,
      $\sigma_{0,\text{phase}}$ from \texttt{filter.measurement\_noise});
\item an elevation down-weighting factor
      $w(\theta^\sat_\rx)^{-1} = (\sin\theta^\sat_\rx)^{-2p_{\text{type}}}$
      with a continuous, type-dependent exponent.  Code and Doppler rows
      use \texttt{filter.code\_elevation\_power} and carrier-phase rows
      \texttt{filter.phase\_elevation\_power}; both default to 1.0, which
      scales the variance by the squared cosecant of the elevation.  When
      the sine of the elevation is at most \texttt{1e-6}, a guard returns
      $\sigma^2_0$ unamplified.  The derivation is in
      \dcref{eq:qc-elevation}{the stochastic-model section of the
      quality-control document};
\item a carrier-to-noise-ratio term $\xi$, parameterised by
      \texttt{snr\_model\_params}.
\end{itemize}
\texttt{filter.stochastic\_model} selects the weighting strategy:
\texttt{constant}, \texttt{elevation} (the default), \texttt{snr} or
\texttt{satellite\_block}.  Only the \texttt{snr} strategy uses the
carrier-to-noise term, and it does not multiply the baseline variance by
a factor: for code and phase rows it replaces the baseline variance by
the variance-additive Hadas--Bosy form built from the observed
carrier-to-noise ratio (coefficients \texttt{code\_a}, \texttt{code\_b},
\texttt{phase\_a}, \texttt{phase\_b}; the ratio is clamped to
20--55~dB-Hz by default; code rows take a per-band scale), then applies
the elevation factor
(\texttt{StochasticModelWeight::computeSnr}).  Cross-correlations between
observations are zero in the undifferenced parameterisation; differencing
introduces the structured correlations treated in \Cref{sec:diff}.
